\chapter*{Аннотация}
\addcontentsline{toc}{chapter}{Аннотация}

Работа посвящена задаче генерации коллективных музыкальных рекомендаций для посетителей заведений и мероприятий методами машинного обучения и совмещает продуктовую проработку сервиса с технической частью на открытом датасете YAMBDA.

В первой части проанализирован рынок музыкального сопровождения коммерческих и событийных пространств, существующие B2B-сервисы и сервисы заказа музыки гостями, а также правовые аспекты публичного воспроизведения и взаимодействия с РАО и ВОИС. Адресуемый рынок постоянных B2B-точек оценён примерно в 25~тыс. объектов, событийный сценарий~--- примерно в 12~тыс. лицензируемых мероприятий в год. На этой основе разработана концепция сервиса <<Резонанс>>: бизнес задаёт рамки музыкального потока, посетители добровольно участвуют в формировании групповой <<волны>>, обработка данных остаётся внутри платформенного контура. Основным механизмом фиксации присутствия выбран BLE-beacon, дополнительным~--- QR-код.

Техническая часть построена как двухуровневая схема: замороженный персональный последовательный скорер (plain SASRec с BCE-функцией потерь) и обучаемые групповые агрегаторы поверх кэша персональных скоров. Скорер воспроизводит индустриальный SASRec-бейзлайн (test~$\mathrm{NDCG@10} = 0{,}0726$ на Listen+ против $0{,}0742$). Для эфемерных групп размера 2--5, синтезированных из YAMBDA, сравниваются ID-based AGREE, GroupIM и две предложенные архитектуры на аудиоэмбеддингах~--- Audio-AGREE и Group Cross-Attention. Аудио-агрегаторы значимо превосходят ID-based бейзлайны по $\mathrm{NDCG@10}$ и $\mathrm{NDCG@20}$; преимущество ярче проявляется на группах среднего размера ($|G| \geq 3$). Дополнительно разработан Android BLE-прототип, иллюстрирующий определение ближайшего сохранённого маяка по силе сигнала.

\section*{Ключевые слова}

Групповые рекомендательные системы; коллективные музыкальные рекомендации; последовательные рекомендации; SASRec; attention-агрегатор; аудиоэмбеддинги; YAMBDA; эфемерные группы; B2B-сервис музыки; BLE-beacon.
